\icode{O-MI-SU-261028}
两杆的水平间距为\(d\)，绳\(AOB\)的长度为\(l\)，物块\(C\)的质量为\(m\)，可得：
\[
    2T\cos\alpha = mg\Rightarrow T = \frac{mg}{2 \cos\alpha } \qquad \sin \alpha = \frac{d}{l}
\]
因此，上下移动系点时，杆间距\(d\)不变时，张力\(T\)不变，如下图所示，动滑轮会建立新的平衡位置

\begin{tikzpicture}[font=\scriptsize]
	\draw (0,0)coordinate(A)node[left]{\(A\)}--(-20:3)--++(70:.1)coordinate(E)--++(-70:.1)coordinate(H)--++(20:3)coordinate(B)node[right]{\(B\)};
	\draw[fill=white] (E) circle(.1);
	\fill (E) circle(.05);
	\draw (E)--($(E)+(0,-.6)$)coordinate(o);
	\draw[fill=orange!50] ($(o)+(-.1,0)$) rectangle($(o)+(.1,-.2)$);
	\draw[dashed,orange] (-20:1.5)--++(70:.1)coordinate(F)--++(-70:.1)--++(20:4.5)coordinate(b)node[right,black]{\(B_1\)};
	\draw[dashed,red] (-20:3)--(-20:4.5)--++(70:.1)coordinate(G)--++(-70:.1)--++(20:1.5)coordinate(C)node[right,black]{\(B_2\)};
	\draw[color=cyan,dashed] (H)--++(-160:{3+.2*tan(20)})coordinate(a)node[left,black]{\(A^{\prime}\)} (a)--(B|-a)coordinate(K)node[right,black]{\(B^{\prime}\)};
	\draw[gray,dashed] (E)--++(0,1.5);
	\draw[gray] ($(E)+(0,{-.1/cos(20)+.3})$) arc(90:160:.3) ($(E)+(0,{-.1/cos(20)+.4})$) arc(90:20:.4);
	\foreach \X in {F,G}{
		\draw[fill=white,draw=blue!50] (\X) circle(.1);
		\fill[color=blue!50] (\X) circle(.05);
    	\draw[dashed,color=blue!50] (\X)--($(\X)+(0,-.6)$)coordinate(O);
	    \draw[fill=blue!50,draw=blue!50] ($(O)+(-.1,0)$) rectangle($(O)+(.1,-.2)$);}
	\draw[fill=orange!60] (0,-3) rectangle(-.1,2) (B|-0,-3)coordinate(I) rectangle($(I)+(.1,5)$);
	\foreach \X in {A,B,b,C,a,K}
		\fill[green!60!blue] (\X) circle(.05);
	\foreach \x in {-.6,-.4,...,6.2}
		\draw[gray] (\x,-3.2)--++(.2,.2);
	\draw[thick] (-.6,-3)--(6.3,-3);
	\node at($(o)+(0,-.2)$)[below]{\(C\)};
	\node at($(E)+(.4,-.1)$){\(O\)};
    \node at ($(E)+(.35,.5)$){\(\alpha\)};
\end{tikzpicture}